% ============================================================================
%  Mini-rapport --- Projet MLOps & Déploiement de modèles
%  Master DEML (M1) --- SupNum --- Prof : Yehdhih ANNA
%  Compilation : pdflatex rapport.tex  (2 passes pour les références)
%
%  >>> CAPTURES D'ECRAN <<<
%  Deposez vos images PNG/JPG dans le dossier  rapport/images/
%  avec EXACTEMENT les noms indiques (s1_mlflow_runs.png, etc.).
%  Tant qu'une image est absente, un cadre gris "Emplacement capture"
%  s'affiche a sa place. Des qu'elle est presente, elle s'insere seule.
% ============================================================================
\documentclass[11pt,a4paper]{article}

\usepackage[utf8]{inputenc}
\usepackage[T1]{fontenc}
\usepackage[french]{babel}
\usepackage{lmodern}
\usepackage[a4paper,margin=2cm]{geometry}
\usepackage{graphicx}
\usepackage{xcolor}
\usepackage{float}
\usepackage{enumitem}
\usepackage{booktabs}
\usepackage{fancyhdr}
\usepackage{titlesec}
\usepackage{tikz}
\usetikzlibrary{shapes.geometric, arrows.meta, positioning, fit, backgrounds}
\usepackage[colorlinks=true, linkcolor=blue, urlcolor=blue, citecolor=blue]{hyperref}

% ---- Couleurs
\definecolor{brand}{HTML}{0F172A}
\definecolor{accent}{HTML}{0369A1}
\definecolor{soft}{HTML}{E2E8F0}
\definecolor{ph}{HTML}{94A3B8}

% ---- En-tete / pied de page
\pagestyle{fancy}
\fancyhf{}
\lhead{\small\textsc{Projet MLOps --- Prédiction du Diabète}}
\rhead{\small SupNum --- DEML M1}
\cfoot{\small\thepage}
\renewcommand{\headrulewidth}{0.4pt}

% ---- Titres de section colorés et compacts
\titleformat{\section}{\large\bfseries\color{accent}}{\thesection.}{0.5em}{}
\titlespacing*{\section}{0pt}{1.1em}{0.5em}
\titleformat{\subsection}{\normalsize\bfseries\color{brand}}{\thesubsection}{0.5em}{}
\titlespacing*{\subsection}{0pt}{0.7em}{0.3em}

\setlist{nosep, leftmargin=1.3em}

% Absorbe les petits débordements de ligne (URLs, mots longs) sans marge cassée.
\setlength{\emergencystretch}{3em}

% ============================================================================
%  Commande de CAPTURE :  \capture{fichier}{legende}{largeur(0-1)}
%  - Si images/<fichier> existe -> l'image est incluse.
%  - Sinon -> cadre gris avec l'instruction (ou prendre la capture).
% ============================================================================
\newcommand{\capture}[3]{%
  \begin{figure}[H]
  \centering
  \IfFileExists{images/#1}{%
    \fbox{\includegraphics[width=#3\linewidth]{images/#1}}%
  }{%
    \setlength{\fboxsep}{0pt}%
    \fcolorbox{ph}{soft!40}{%
      \begin{minipage}[c][3.6cm][c]{#3\linewidth}
        \centering\color{ph}\itshape
        \textbf{[ Emplacement capture --- \texttt{\detokenize{images/#1}} ]}\\[6pt]
        #2
      \end{minipage}}%
  }
  \caption{#2}
  \end{figure}
}

% ============================================================================
\begin{document}

% ------------------------------------------------------------------ TITRE
\begin{center}
  {\LARGE\bfseries\color{brand} Prédiction du Diabète --- Pipeline MLOps de bout en bout}\\[4pt]
  {\large Module : MLOps et Déploiement de modèles}\\[2pt]
  {\normalsize Institut Supérieur du Numérique --- SupNum \quad$\cdot$\quad Master DEML (M1) --- 2026}\\[6pt]
  {\normalsize Enseignant : Yehdhih ANNA \quad$\cdot$\quad Date : 3 juillet 2026}\\[8pt]
  \begin{tabular}{l}
    \textbf{Groupe 1 :} Bechir Mady (22014) \quad$\cdot$\quad Abdarahman Abdarahman (22053) \quad$\cdot$\quad Aly Wedad (22086)
  \end{tabular}
\end{center}
\vspace{0.4em}
\hrule
\vspace{0.8em}

% ------------------------------------------------------------------ LIENS
\section{Liens d'accès}
\begin{itemize}
  \item \textbf{Repository GitHub :} \url{https://github.com/AymanMady/Projet-MLOps}
  \item \textbf{Serveur MLflow (Tracking + Registry, EC2 \#1) :} \url{http://54.95.120.104:5000}
  \item \textbf{Application Flask (EC2 \#2) :} \url{http://54.95.120.104:5001} \quad{\small\color{ph}}
  \item \textbf{Stockage S3 :} bucket \texttt{m1-2026-mlops-project-group1} (données brutes / prétraitées, cache DVC, artefacts).
\end{itemize}

% ------------------------------------------------------------------ PROBLEME ML
\section{Présentation du problème ML et du dataset}
\subsection*{Problème traité}
Il s'agit d'un problème de \textbf{classification binaire supervisée} en santé : prédire si un
patient est \emph{diabétique} (\texttt{Outcome = 1}) ou \emph{non diabétique} (\texttt{Outcome = 0})
à partir de mesures cliniques simples. L'objectif applicatif est un outil d'aide au dépistage
précoce, exposé via une interface web où un utilisateur saisit les mesures d'un patient et obtient
une prédiction immédiate.

\subsection*{Jeu de données}
Nous utilisons le dataset \textbf{Pima Indians Diabetes} (\textasciitilde768 observations, 8 variables
explicatives numériques + 1 cible). Les données brutes sont versionnées avec DVC et stockées sur S3.

\begin{table}[H]
\centering\small
\begin{tabular}{@{}lll@{}}
\toprule
\textbf{Variable} & \textbf{Description} & \textbf{Type} \\ \midrule
Pregnancies & Nombre de grossesses & entier \\
Glucose & Glycémie (test de tolérance) & numérique \\
BloodPressure & Pression artérielle diastolique (mm Hg) & numérique \\
SkinThickness & Épaisseur du pli cutané tricipital (mm) & numérique \\
Insulin & Insuline sérique (mu U/ml) & numérique \\
BMI & Indice de masse corporelle & numérique \\
DiabetesPedigreeFunction & Antécédents familiaux (score) & numérique \\
Age & Âge (années) & entier \\ \midrule
\textbf{Outcome} & \textbf{Cible : 1 = diabétique, 0 = non} & \textbf{binaire} \\ \bottomrule
\end{tabular}
\caption{Variables du dataset Pima Indians Diabetes.}
\end{table}

\noindent\textbf{Modèle \& métriques.} Nous entraînons un \texttt{RandomForestClassifier}
(scikit-learn ; \texttt{n\_estimators=200}, \texttt{max\_depth=8}). Les données sont nettoyées
(suppression des doublons, imputation médiane) puis mises à l'échelle (\texttt{StandardScaler}).
Les performances sont suivies via l'\textbf{accuracy}, le \textbf{F1-score}, la \textbf{précision}
et le \textbf{rappel}, tous logués dans MLflow.

% ------------------------------------------------------------------ ARCHITECTURE
\section{Architecture technique mise en place}

\begin{figure}[H]
\centering
\begin{tikzpicture}[
  font=\small,
  box/.style={rectangle, rounded corners, draw=accent, thick, fill=soft!45,
              align=center, minimum height=1.5cm, minimum width=3.2cm, inner sep=5pt},
  s3/.style={cylinder, shape border rotate=90, aspect=0.25, draw=accent, thick,
             fill=accent!12, align=center, minimum height=2cm, minimum width=3cm},
  arr/.style={-{Stealth[length=3mm]}, thick, accent},
  lbl/.style={font=\scriptsize\itshape, text=brand}
]
  % Noeuds
  \node[box] (pc) {\textbf{PC local}\\ scripts \texttt{src/}\\ DVC + \texttt{params.yaml}};
  \node[s3, right=2.6cm of pc] (s3) {\textbf{Amazon S3}\\ \scriptsize raw / processed\\ \scriptsize dvcstore + artefacts};
  \node[box, above right=0.3cm and 2.6cm of s3] (ec1) {\textbf{EC2 \#1 --- MLflow}\\ Tracking (:5000)\\ Model Registry};
  \node[box, below right=0.3cm and 2.6cm of s3] (ec2) {\textbf{EC2 \#2 --- Flask}\\ App web (:5001)\\ formulaire + prédiction};

  % Fleches
  \draw[arr] (pc) -- node[lbl, above]{données} node[lbl,below]{DVC push/pull} (s3);
  \draw[arr] (s3) -- node[lbl, above, sloped]{artefacts} (ec1);
  \draw[arr] (pc.north) to[out=60,in=180] node[lbl, above]{log params / métriques / modèle} (ec1.west);
  \draw[arr] (ec1) -- node[lbl, right, align=left]{modèle\\ \og champion \fg} (ec2);
  \draw[arr] (s3) -- node[lbl, below, sloped]{secours} (ec2);
\end{tikzpicture}
\caption{Architecture : le poste local orchestre le pipeline, S3 centralise les données/artefacts,
EC2\#1 héberge MLflow (tracking + registry) et EC2\#2 sert le modèle via Flask.}
\end{figure}

\noindent\textbf{Rôle de chaque composant :}
\begin{itemize}
  \item \textbf{PC local} --- exécute les scripts (\texttt{preprocess} $\rightarrow$ \texttt{train}
        $\rightarrow$ \texttt{evaluate}), pilote DVC et logue vers MLflow.
  \item \textbf{Amazon S3} (\texttt{m1-2026-mlops-project-group1}) --- dépôt partagé : données brutes,
        données prétraitées, \emph{remote} DVC (\texttt{dvcstore}) et artefacts/modèles MLflow.
  \item \textbf{EC2 \#1 --- MLflow} --- serveur de \emph{tracking} des expériences et \emph{Model Registry}
        (gestion des versions + alias \og champion \fg).
  \item \textbf{EC2 \#2 --- Flask} --- application web qui charge le modèle \og champion \fg{} du
        Model Registry et renvoie la prédiction ; repli sur le modèle local si le Registry est
        indisponible.
\end{itemize}

% ------------------------------------------------------------------ PIPELINE ML + MLFLOW
\section{Pipeline ML et suivi des expériences (MLflow)}
Trois scripts, intégrés à MLflow dès le départ :
\begin{enumerate}
  \item \texttt{src/preprocess.py} --- nettoyage, imputation médiane, \texttt{StandardScaler} ;
        écrit les données prétraitées sur S3.
  \item \texttt{src/train.py} --- entraîne le RandomForest, logue les \textbf{hyperparamètres} et les
        \textbf{métriques} (accuracy, F1), et enregistre le modèle comme artefact.
  \item \texttt{src/evaluate.py} --- évalue (accuracy, F1, précision, rappel) et \textbf{n'enregistre
        que le meilleur modèle} dans le Registry (si accuracy $\geq$ seuil \emph{et} $\geq$ meilleure
        version existante), puis lui affecte l'alias \og champion \fg.
\end{enumerate}

\capture{s1_mlflow_runs.png}{MLflow --- liste des runs de l'expérience \og Student\_Experiment \fg{} (preprocess, train, evaluate).}{0.95}

\capture{s2_mlflow_run_detail.png}{MLflow --- détail d'un run : hyperparamètres (n\_estimators, max\_depth\dots{}) et métriques (accuracy, F1, précision, rappel).}{0.95}

\capture{s3_mlflow_registry.png}{MLflow --- Model Registry : modèle \og diabetes-classifier \fg{}, ses versions et l'alias \og champion \fg{}.}{0.95}

% ------------------------------------------------------------------ DVC + S3
\section{Données \& reproductibilité (DVC + S3)}
Les données sont versionnées avec \textbf{DVC}, en utilisant S3 comme \emph{remote storage}
(dossier \texttt{dvcstore} du bucket).

\capture{s4_dvc_terminal.png}{DVC --- terminal montrant \texttt{dvc status} / \texttt{dvc push} (ou \texttt{dvc dag}) et le versioning des données.}{0.9}

\capture{s5_s3_bucket.png}{Console AWS S3 --- bucket \og m1-2026-mlops-project-group1 \fg{} avec les dossiers data/raw, data/processed et dvcstore.}{0.9}



% ------------------------------------------------------------------ DEPLOIEMENT
\section{Déploiement --- Application Flask}
L'application \texttt{app/app.py} expose un formulaire (les 8 features cliniques). À la soumission,
elle charge le modèle \og champion \fg\ depuis le Model Registry MLflow, effectue la prédiction et
affiche le résultat (\og Diabétique \fg{} / \og Non diabétique \fg{}). Elle est déployée sur une seconde
instance EC2 (port 5001).

\begin{minipage}[t]{0.47\linewidth}
\capture{s7_flask_form.png}{Interface Flask --- formulaire de saisie des mesures.}{1.0}
\end{minipage}\hfill\begin{minipage}[t]{0.47\linewidth}
\capture{s8_flask_result.png}{Interface Flask --- résultat de la prédiction affiché.}{1.0}
\end{minipage}

% ------------------------------------------------------------------ GIT
\section{Collaboration \& historique Git}
Le dépôt suit un modèle multi-branches : chaque membre dispose d'une branche personnelle
(\texttt{bechir-22014}, \texttt{abdrahman}, \texttt{aliy-22086}) fusionnée dans \texttt{main} via des
\emph{Pull Requests}, produisant un arbre Git non linéaire avec plusieurs \emph{merges}.

\capture{s9_git_graph.png}{Graphe Git (GitHub Network ou \og git log --graph \fg{}) montrant les branches par membre et les merges.}{0.9}

\vfill
\begin{center}\small\color{ph}
Rapport projet MLOps --- Groupe 1 --- SupNum DEML M1.
\end{center}

\end{document}
